\section{Related Inverse Problems and Algorithmic Context}\label{sec:related}
\subsection{Positive regularization and multidimensional structure}
Positive Laplace inversion belongs to the broader class of first-kind Fredholm equations. Classical exponential-decay analysis already distinguishes recovering a decay representation from resolving individual exponentials \cite{discrete}. Small singular values make unregularized solutions sensitive to perturbations, so numerical rank, discretization and regularization must be considered together \cite{hansen}. The nonnegative smoothing formulation of Butler, Reeds and Dawson \cite{brd}, CONTIN \cite{contin}, and uniform-penalty inversion (UPEN) \cite{upen} provide established alternatives for stabilization. Their different penalties express different preferences for admissible solutions; none makes a finite data vector identify an unrestricted positive distribution.

Efficient optimization does not remove that information limit. Adaptive truncation and the FLINT implementation address the cost of estimating NMR relaxation distributions \cite{flint}, while accelerated proximal methods such as FISTA provide a general optimization framework for composite convex objectives \cite{fista}. Their convergence statements concern the specified objectives under their assumptions. They do not automatically apply to the bilinear factorization, data-dependent screening, or changing component order used here.

Methods that exploit tensor structure in Fredholm inversion \cite{separable2d}, UPEN in two dimensions \cite{upen2d}, and uniform multipenalty regularization \cite{mupen} provide precedents for exploiting operator structure and adapting regularization. A DOSY matrix has a particular asymmetry: diffusion enters a Laplace kernel, whereas the chemical-shift dimension indexes spectra. Joint processing of chemical shifts is therefore not, by itself, a second Laplace inversion. Nor can performance under one regularization rule establish the adequacy of another rule's smoothing constants. Our frozen MF example illustrates one declared discretization and pair of constants; selecting them adaptively over a wider application domain remains separate work.

\subsection{Spectral unmixing, finite exponentials and identification}
Multivariate curve resolution \cite{mcr} and nonnegative matrix factorization \cite{nmf} establish the general idea of explaining observations through shared nonnegative factors. In TRAIn-MF the attenuation factor is constrained further: it must have the form $KS$, where $K$ is a known diffusion kernel. Column normalization of $S$ removes a scale freedom but does not remove all factorization ambiguities. A single shared profile may be broad or multimodal, so its factor count cannot be read as a molecular count.

Within DOSY, SCORE \cite{score}, OUTSCORE \cite{outscore}, and LOCODOSY \cite{locodosy} are direct precedents for component resolution using information across resonances. Global sharing and local covariance order offer different responses to spectral overlap: a component need not contribute at every chemical shift, and a local window need not have the mixture's global order. These ideas are central context for the present support matrix, rather than evidence that spectral correlation is a new concept. The GNAT processing environment \cite{gnat} also demonstrates the importance of distinguishing an inversion routine from the surrounding NMR processing workflow. Here phase, baseline and physical diffusion weighting must be supplied correctly before an inversion is interpreted.

Finite exponential estimation has additional tools, including ESPIRA, which uses rational approximation and structured matrices to estimate exponential parameters \cite{espira}. Its sampling and model assumptions must be checked before transferring it to a positive DOSY analysis. More generally, identifiability is a property of a specified model class \cite{teicher}. Uniqueness in a bounded-order noiseless atomic class is compatible with nonuniqueness among unrestricted positive measures, and neither statement guarantees stable recovery of nearby rates in noise. Our propositions make those distinctions explicit; the comparison does not claim to supersede these established theories.

\subsection{Learned reconstructions and the scope of the comparison}
DRECT introduces deep learning for Laplace NMR reconstruction \cite{drect}, and DREAM adds uncertainty-informed reconstruction \cite{dream}. Learned estimators can encode information through their training distributions as well as through the physical kernel. Consequently, assessing them requires declared training priors, independent test cases, and checks under distribution shift. The present selector's conditional Gaussian calculations and numerical flags do not provide an equivalent learned uncertainty model or a calibrated chemical confidence interval.

These references extend the methodological context, not the experimental comparison. No new performance ranking of FLINT, UPEN, SCORE, ESPIRA, DRECT or DREAM is inferred from the literature. The results section retains the same frozen estimators, data and metrics. The specific contribution assessed here is the combination of a declared positive representation, support-aware conditional fitting and predictive selection, with explicit limitations on identifiability and inference.
